\documentclass[11pt,a4paper]{article}

\usepackage[utf8]{inputenc}
\usepackage[T1]{fontenc}
\usepackage{microtype}
\usepackage[margin=2.5cm]{geometry}
\usepackage{amsmath,amssymb}
\usepackage{graphicx}
\usepackage{booktabs}
\usepackage{longtable}
\usepackage{array}
\usepackage{caption}
\captionsetup[table]{position=top,skip=6pt}

\usepackage[backend=biber,style=numeric-comp,sorting=none]{biblatex}
\addbibresource{../../bibliography/clifford-seminar-group.bib}

\usepackage[colorlinks=true]{hyperref}
\usepackage{cleveref}

\newcommand{\Cl}{\mathcal{C}\ell}
\newcommand{\R}{\mathbb{R}}

\title{\LARGE\bfseries Clifford Algebras in Theoretical Physics\\[0.5em]
\large Seminar Methodology, Modular Syllabus, and Thematic Progression}
\author{Clifford Seminar Group\\
  \small Leader: Julian L. Avila-Martinez \\
  \small Universidad Distrital Francisco Jos{\'e} de Caldas}
\date{}

\begin{document}

\maketitle

\section{Overview}
\label{sec:overview}

This seminar investigates invariant geometric formulations of theoretical
physics through Clifford algebras.
The Kepler problem and the associated Coulomb potential serve as the central
dynamical backbone and comparative benchmark across classical mechanics, spinor
regularization, relativistic spacetime dynamics, and relativistic quantum
mechanics.
While this central force problem guides the thematic progression, the curriculum
is designed to accommodate systematic deviations into adjacent physical and
mathematical domains, including coordinate-free differential forms, monogenic
analysis, gauge theories of gravitation, and geometric thermodynamics.

The colloquium operates as a collaborative working seminar focused on direct
blackboard derivations and critical appraisal of physical foundations.
The syllabus is structured into sequential modular topics.
A given topic may span one or multiple 90-minute sessions depending on the
analytical depth of the derivations and the scope of plenary debate.
This modular design preserves structural continuity when topics expand or when
tangential inquiries are pursued.

Exposition duties are distributed across senior--junior tandems to foster
rigorous peer mentorship.

\section{Methodology}
\label{sec:methodology}

\subsection{Working Roles}
\label{subsec:roles}

The seminar must avoid passive lecture formats.
Each session relies on four distinct responsibilities.

\begin{description}
  \item[Referat (Lead Exposition)]
  Delivered jointly by the assigned tandem.
  The presentation focuses on direct, blackboard derivations of primary
  equations.
  Conceptual assumptions and algebraic transitions take precedence over
  projected slides.
  \item[Korreferat (Respondent Critique)]
  An assigned participant presents a critical commentary immediately following
  the Referat.
  The respondent examines implicit hypotheses, compares the results with
  traditional coordinate or matrix formalisms, identifies potential pathologies,
  and formulates targeted questions for discussion.
  \item[Thesenpapier (Position Paper)]
  A concise memorandum circulated by the presenting tandem at least five days
  prior to the session.
  It contains formal definitions, key algebraic identity statements, and
  explicit theses concerning physical interpretation.
  Participants study this document in advance.
  \item[Sitzungsprotokoll (Session Record)]
  A rotating participant documents the derivation steps, the resolution of
  theoretical debates, and open problems.
  The finalized protocol is archived in the shared repository within seven days
  following the session.
\end{description}

\subsection{Tandem Collaboration}
\label{subsec:tandems}

Beginning with Topic~2, sessions pair an advanced participant (senior) with an
earlier-stage student (junior).
The senior ensures mathematical coherence and contextualizes advanced algebraic
techniques.
The junior works directly through the derivation machinery, receiving structured
guidance in multilinear structures and relativistic dynamics.
Responsibilities for blackboard presentation are shared equally between both
speakers.

\subsection{Session Structure}
\label{subsec:structure}

Preparatory meetings with the seminar leader occur two weeks and five days prior
to each session.
Every attendee must submit at least one written question on the Thesenpapier 24
hours before meeting.
The 90-minute session adheres to the allocation in \cref{tab:session_structure}.

\begin{table}[ht]
\centering
\caption{Allocation of time per 90-minute seminar session.}
\label{tab:session_structure}
\begin{tabular}{@{}llll@{}}
\toprule
Phase & Duration & Content & Responsible \\
\midrule
I   & 40 min & Live blackboard derivation (Referat) & Presenting tandem \\
II  & 10 min & Critical response and opening theses (Korreferat) & Respondent \\
III & 35 min & Plenary debate and interpretation & All participants \\
IV  & 5 min  & Synthesis of open questions and preview & Seminar leader \\
\bottomrule
\end{tabular}
\end{table}

\section{Schedule and Syllabus}
\label{sec:schedule}

Primary reference works are designated by short abbreviations.
VR indicates Vaz and da Rocha, \emph{An Introduction to Clifford Algebras and
Spinors} \cite{vazIntroductionCliffordAlgebras2016}; DL indicates Doran and
Lasenby, \emph{Geometric Algebra for Physicists}
\cite{doranGeometricAlgebraPhysicists2013}; HS indicates Hestenes and Sobczyk,
\emph{Clifford Algebra to Geometric Calculus}
\cite{hestenesCliffordAlgebraGeometric1984}.
Specific chapters and foundational literature are cited directly in
\cref{tab:syllabus}.

\renewcommand{\arraystretch}{1.28}
\begin{longtable}{%
  >{\raggedright\arraybackslash}p{0.8cm}%
  >{\raggedright\arraybackslash}p{3.5cm}%
  >{\raggedright\arraybackslash}p{6.3cm}%
  >{\raggedright\arraybackslash}p{4.4cm}%
}
\caption{Modular seminar syllabus and thematic sequence.}
\label{tab:syllabus}\\
\toprule
Topic & Subject & Mathematical and Physical Content & Literature \\
\midrule
\endhead

\multicolumn{4}{@{}l}{%
  \textit{Part A: Foundations and Geometric Analysis}%
}\\
1 & Dual spaces, multilinear forms, and quadratic spaces
  & Dual linear space $V^*$, dual pairing $\langle\cdot,\cdot\rangle$, tensor
    algebra $T(V)$, symmetric bilinear forms, quadratic forms $q(v)$,
    polarization, non-degeneracy, and signature $(p,q)$ over $\R$.
  & VR \cite[Chap.~1]{vazIntroductionCliffordAlgebras2016};
    Axler \cite[Chaps.~1--2]{axlerLinearAlgebraDone2015};
    Fraleigh \cite[Chaps.~1.1--1.2, 7.33]{fraleighFirstCourseAbstract2021}.
  \\[3pt]
2 & Universal Clifford construction and involutions
  & Ideal $\mathcal{I}_q=\langle v\otimes v-q(v)1\rangle$, quotient algebra
    $\Cl(V,q)$, universal property, $\mathbb{Z}_2$-grading, grade involution
    $\hat{u}$, reversion $\tilde{u}$, Clifford conjugation $\bar{u}$,
    geometric product $uv=u\cdot v+u\wedge v$.
  & \\[3pt]
3 & Multivector structures, Lie algebras, and Spin groups
  & Module isomorphism $\Cl(V,q)\cong\bigwedge V$, blade projections,
    bivector Lie algebra $\Cl^2(V,q)\cong\mathfrak{so}(p,q)$, Clifford and
    Lipschitz groups, $\mathrm{Pin}(p,q)$, $\mathrm{Spin}(p,q)$, and rotor
    transformations $v\mapsto Rv\tilde{R}$.
  & \\[3pt]
4 & Clifford analysis and monogenic fields
  & Vector derivative $\nabla$, Cauchy--Riemann generalized operator, monogenic
    multivector fields $\nabla\Phi=0$, Cauchy kernel, Borel--Pompeiu integral
    formula, and coordinate-free multipole expansions of extended Keplerian
    sources.
  & \\
\midrule

\multicolumn{4}{@{}l}{%
  \textit{Part B: Classical Kepler Dynamics in $\Cl(3,0)$}%
}\\
5 & Central forces and orbital angular momentum
  & Equation of motion $\dot{p}=-(k/r^3)r$, orbital angular momentum bivector
    $L=r\wedge p$, conservation proof $\dot{L}=0$, planar confinement
    $r\wedge L=0$, coordinate-free duality via volume pseudoscalar $I_3$.
  & \\[3pt]
6 & Laplace--Runge--Lenz vector and conics
  & Multivector construction $A=\frac{1}{mk}pL-\hat{r}$, algebraic proof of
    $\dot{A}=0$ without coordinate frames, orbit derivation via scalar
    contraction $r\cdot A$, and invariant classification of conics.
  & \\
\midrule

\multicolumn{4}{@{}l}{%
  \textit{Part C: Regularization, Harmonic Analysis, and Hidden Symmetries}%
}\\
7 & Kustaanheimo--Stiefel transformation
  & Regularization of collision singularities at $r\to 0$, spinor mapping
    $x=u\sigma_3\tilde{u}$ with $u\in\Cl^+(3,0)\cong\mathbb{H}$, Sundman time
    rescaling $dt=r\,ds$, and $U(1)$ gauge freedom.
  & \\[3pt]
8 & Four-dimensional oscillator and $\mathrm{SO}(4)$
  & Linearization to $u''+\frac{|E|}{2}u=0$, four-dimensional isotropic harmonic
    oscillator, realization of $\mathfrak{so}(4)\cong\mathfrak{su}(2)\oplus
    \mathfrak{su}(2)$ dynamical symmetry via $L$ and $A$, bound-state
    degeneracy.
  & \\[3pt]
9 & Clifford--Fourier analysis and Fock projection
  & Clifford--Fourier transform with pseudoscalar kernel $e^{-I_3 k\cdot x}$,
    multivector uncertainty relations, Fock stereographic projection of the
    Coulomb problem onto $S^3$, and bound states as spherical monogenics.
  & \\
\midrule

\multicolumn{4}{@{}l}{%
  \textit{Part D: Relativistic Dynamics in Spacetime Algebra $\Cl(1,3)$}%
}\\
10 & Spacetime algebra and observer split
  & Dirac generators $\{\gamma_\mu\}$, signature $(+,-,-,-)$, observer split
    via $\gamma_0$, relative spatial bivectors $\sigma_k=\gamma_k\gamma_0$,
    proper Lorentz rotors in $\mathrm{Spin}^+(1,3)$, rotor kinematics
    $\dot{R}=\frac{1}{2}\Omega R$, Thomas precession.
  & \\[3pt]
11 & Relativistic electrodynamics and Lorentz force
  & Spacetime vector derivative $\nabla=\gamma^\mu\partial_\mu$, Faraday
    bivector $F=E+I_4B$, Maxwell system $\nabla F=J$, relativistic equation of
    motion $\dot{p}=eF\cdot u$, proper-time rotor integration.
  & \\[3pt]
12 & Relativistic central forces and perihelion advance
  & Relativistic central potentials, breakdown of LRL vector conservation,
    energy-dependent dynamical corrections, and calculation of
    special-relativistic perihelion advance.
  & \\[3pt]
13 & Gauge Theory Gravity and planetary orbits
  & Position gauge field $\bar{h}(a)$ and rotation gauge field
    $\bar{\Omega}(a)$ on Minkowski background, spherically symmetric
    solutions, geodesic rotor trajectories, and general-relativistic orbital
    precession.
  & \\[3pt]
14 & Geometric thermodynamics in spacetime algebra
  & Conserved particle flux $N=nu$, stress-energy operator $T(a)$, entropy
    current $S=su$, invariant thermal 1-vector $\beta=u/(k_{\mathrm{B}}T)$,
    Killing equilibrium $\nabla\wedge\beta=0$, and Tolman--Ehrenfest thermal
    gradients in static gravitational fields.
  & \\
\midrule

\multicolumn{4}{@{}l}{%
  \textit{Part E: Quantum Kepler Problem and Dirac--Hestenes Theory}%
}\\
15 & Real Pauli algebra and spatial spin
  & Spatial algebra $\Cl(3,0)$ without matrix representations, unit spatial
    vectors as spin observables, Pauli multivector wave equation, spin-orbit
    coupling from the geometric product.
  & \\[3pt]
16 & Dirac--Hestenes field equation
  & Spacetime algebra representation of the Dirac equation, even multivector
    field $\Psi\in\Cl^+(1,3)$, spin plane $\gamma_2\gamma_1$, polar form
    $\Psi=\rho^{1/2}e^{I_4\beta/2}R$, probability current
    $J=\Psi\gamma_0\tilde{\Psi}$, spin bivector $S=\frac{1}{2}\hbar\Psi\gamma_3
    \tilde{\Psi}$.
  & \\[3pt]
17 & Exact solution of the relativistic hydrogen atom
  & Separation of variables in spherical coordinates within $\Cl^+(1,3)$, rotor
    spherical harmonics, radial differential equations in the even sub-algebra,
    derivation of the Sommerfeld--Dirac energy spectrum.
  & \\[3pt]
18 & Fine structure, Darwin term, and zitterbewegung
  & Non-relativistic expansion of the Dirac--Hestenes equation, geometric
    origin of the Darwin correction, rotor phase dynamics, lightlike field
    circulation, field-theoretic interpretation of zitterbewegung.
  & \\
\bottomrule
\end{longtable}

\section{Tandem Assignments}
\label{sec:tandems_table}

The assignments for lead expositions, junior collaborators, and designated
respondents are compiled in \cref{tab:tandems}.

\begin{table}[ht]
\centering
\caption{Register of seminar leaders, tandems, and respondents across topics.}
\label{tab:tandems}
\renewcommand{\arraystretch}{1.15}
\begin{tabular}{@{}cp{4.6cm}p{3.3cm}p{3.3cm}p{2.6cm}@{}}
\toprule
Topic & Subject & Senior / Lead & Junior & Respondent \\
\midrule
1  & Dual spaces and quadratic forms   & J.~L. Avila-Martinez & --- & \\[2pt]
2  & Universal Clifford construction  & J.~L. Avila-Martinez &
     C.~Martinez & \\[2pt]
3  & Multivector structures and Spin  & & & \\[2pt]
4  & Clifford analysis and monogenics & & & \\[2pt]
5  & Central forces and bivectors     & & & \\[2pt]
6  & LRL vector and conics            & & & \\[2pt]
7  & KS transformation                & & & \\[2pt]
8  & 4D oscillator and $\mathrm{SO}(4)$ & & & \\[2pt]
9  & Clifford--Fourier and Fock projection & & & \\[2pt]
10 & Spacetime algebra                & & & \\[2pt]
11 & Electrodynamics in STA           & & & \\[2pt]
12 & Relativistic central forces      & & & \\[2pt]
13 & Gauge Theory Gravity             & & & \\[2pt]
14 & Geometric thermodynamics         & & & \\[2pt]
15 & Pauli algebra and spin           & & & \\[2pt]
16 & Dirac--Hestenes equation         & & & \\[2pt]
17 & Relativistic hydrogen            & & & \\[2pt]
18 & Fine structure and zitterbewegung & & & \\
\bottomrule
\end{tabular}
\end{table}

\printbibliography

\end{document}
